\subsection{Architecture discovery protocol and resources}\label{app:architecture_search}

Each request supplies the original twelve labeled examples or nine integer indices, input/output shapes and the deployment domain. It supplies no task name, formula, initial weights or prescribed architecture. Extrapolation is the primary objective. Resource minimization is secondary. There are three independent first responses per task, and no response is selected using test data. Every graph is frozen before evaluation. The two invalid responses remain failed trials.

Allowed operators are affine maps, rectified linear units (ReLU), hyperbolic tangent, sigmoid, elementwise addition and multiplication, static tensor indexing and concatenation, reshape, transpose, matrix multiplication, softmax and one-dimensional convolution. Optional recurrence repeats a shared cell. This library permits multilayer perceptrons, recurrent neural networks, convolutional neural networks and attention networks, as well as multiplicatively gated arithmetic graphs. It bounds every proposal at 64 nodes and 4,096 scalar coefficients. A tanh node directly implements hyperbolic tangent, and a multiplication node directly implements square. Both operations therefore affect capacity and the appropriate interpretation of comparisons with the fixed ReLU architectures.

Each floating-point operation (FLOP) is a scalar addition, multiplication or division. Dot products count every multiply and accumulation, including fixed zeros and ones. Comparisons and transcendental evaluations are listed separately. The coefficient count includes every stored scalar constant, bias, matrix entry and initial-state value. Each coefficient takes eight bytes. Peak workspace counts live binary64 input, state, intermediate and output buffers under the interpreter\textquotesingle s dependency-based lifetime rule. Graph outputs are materialized. Python, allocator and matrix-library workspaces are excluded. Serialized executable JavaScript Object Notation (JSON) reports the full graph-file size, including topology and numeric literals. This serialization metric is separate from runtime memory. For recurrence, costs use one example starting from the initial state; the listed worst case is index 64 or a six-element stream. A comparison of costs therefore requires the same operations, precision, input length and memory accounting.

Zero stored coefficients is the elementary lower bound on a nonnegative coefficient count. The square, tanh and addition proposals attain that bound in this operator library. The square graph also uses one arithmetic multiplication. A nonconstant scalar output requires at least one stored output scalar under the materialized-output convention, or eight bytes in binary64. The observed 16-byte square workspace includes input and output. These simple bounds provide no certificate that the other proposed architectures minimize coefficients, workspace or arithmetic.

The three sine constructions illustrate competing resource objectives. Their coefficient counts were 19, 14 and 10, while peak live workspaces were 32, 40 and 88 bytes and arithmetic counts were 32, 27 and 27. All three produced the same measured outer RMSE. Reducing coefficient storage therefore did not minimize workspace in this sample.

The independently drawn confirmation set uses a post-freeze random seed. It contains 2,048 new central and 2,048 new outer inputs per stochastic task. The integer tasks have fixed held-out indices 17--32 and 33--64. All fifteen exact stochastic-task proposals also matched every independently drawn confirmation output. The six exact integer-sequence proposals concern fixed indices. All three running-sum proposals also matched all 2,048 length-32 central streams. The complete scores are in the data archive.

\begin{table}[p]\centering\footnotesize

\caption{Every architecture proposal on the original outer test. Exact is the number of whole examples matching the reference. FLOPs, comparisons, transcendental calls and workspace use the largest input in the split. Dashes indicate a rejected executable graph. Weight bytes count stored coefficients once, including recurrent initial state. Graph bytes measure the serialized graph file, including numeric literals. Live bytes measure peak live tensor storage; all three sizes are in bytes.}\label{tab:architecture_all}

\setlength{\tabcolsep}{2.3pt}

\begin{tabular}{@{}lr r r r r r r r r@{}}\toprule

Target & Seed & Exact & RMSE & \shortstack{Weight\\bytes} & \shortstack{Graph\\bytes} & FLOPs & \shortstack{Compar-\\isons} & \shortstack{Function\\calls} & \shortstack{Live\\bytes}\\\midrule

Square & 164 & 4096/4096 & 0 & 0 & 109 & 1 & 0 & 0 & 16\\

Square & 317 & 4096/4096 & 0 & 0 & 109 & 1 & 0 & 0 & 16\\

Square & 583 & 4096/4096 & 0 & 0 & 109 & 1 & 0 & 0 & 16\\

\addlinespace[2pt]

Sine & 164 & 273/4096 & $7.79\times10^{-16}$ & 152 & 1580 & 32 & 0 & 0 & 32\\

Sine & 317 & 273/4096 & $7.79\times10^{-16}$ & 112 & 1688 & 27 & 0 & 0 & 40\\

Sine & 583 & 273/4096 & $7.79\times10^{-16}$ & 80 & 1937 & 27 & 0 & 0 & 88\\

\addlinespace[2pt]

Positive self-power & 164 & 0/2048 & $1.68\times10^{4}$ & 744 & 4926 & 250 & 0 & 0 & 1040\\

Positive self-power & 317 & invalid & $\infty$ & -- & -- & -- & -- & -- & --\\

Positive self-power & 583 & invalid & $\infty$ & -- & -- & -- & -- & -- & --\\

\addlinespace[2pt]

Hyperbolic tangent & 164 & 4096/4096 & 0 & 0 & 101 & 0 & 0 & 1 & 16\\

Hyperbolic tangent & 317 & 4096/4096 & 0 & 0 & 81 & 0 & 0 & 1 & 16\\

Hyperbolic tangent & 583 & 4096/4096 & 0 & 0 & 101 & 0 & 0 & 1 & 16\\

\addlinespace[2pt]

Addition & 164 & 2048/2048 & 0 & 0 & 202 & 1 & 0 & 0 & 32\\

Addition & 317 & 2048/2048 & 0 & 0 & 220 & 1 & 0 & 0 & 32\\

Addition & 583 & 2048/2048 & 0 & 0 & 220 & 1 & 0 & 0 & 32\\

\addlinespace[2pt]

Subtraction & 164 & 2048/2048 & 0 & 16 & 119 & 3 & 0 & 0 & 24\\

Subtraction & 317 & 2048/2048 & 0 & 16 & 119 & 3 & 0 & 0 & 24\\

Subtraction & 583 & 2048/2048 & 0 & 16 & 119 & 3 & 0 & 0 & 24\\

\addlinespace[2pt]

Running sum & 164 & 2048/2048 & 0 & 8 & 173 & 6 & 0 & 0 & 112\\

Running sum & 317 & 2048/2048 & 0 & 8 & 173 & 6 & 0 & 0 & 112\\

Running sum & 583 & 2048/2048 & 0 & 8 & 173 & 6 & 0 & 0 & 112\\

\addlinespace[2pt]

Sort four & 164 & 1096/2048 & $5.23\times10^{-16}$ & 8 & 1546 & 25 & 5 & 0 & 96\\

Sort four & 317 & 1128/2048 & $5.15\times10^{-16}$ & 64 & 1484 & 35 & 5 & 0 & 112\\

Sort four & 583 & 1096/2048 & $5.23\times10^{-16}$ & 8 & 1713 & 25 & 5 & 0 & 88\\

\addlinespace[2pt]

Fibonacci & 164 & 32/32 & 0 & 16 & 389 & 64 & 0 & 0 & 48\\

Fibonacci & 317 & 32/32 & 0 & 16 & 389 & 64 & 0 & 0 & 48\\

Fibonacci & 583 & 32/32 & 0 & 16 & 389 & 64 & 0 & 0 & 48\\

\addlinespace[2pt]

Triangular & 164 & 32/32 & 0 & 8 & 243 & 3 & 0 & 0 & 24\\

Triangular & 317 & 32/32 & 0 & 16 & 205 & 3 & 0 & 0 & 24\\

Triangular & 583 & 32/32 & 0 & 16 & 205 & 3 & 0 & 0 & 24\\

\addlinespace[2pt]

\bottomrule\end{tabular}\end{table}

\clearpage

\subsection{Complete architecture-selection prompts and final answers}\label{app:architecture_prompts}

The following two complete prompts use the first data seed, 164, for square and sine. They are literal request texts, with line wrapping added for printing. The archive supplies all 30 prompts and replies. The network language and objective were fixed before dispatch. Each response contained a public construction summary and a proposed mathematical certificate. These are model-supplied final-answer explanations; no hidden internal reasoning trace was recorded. Numerical outcomes and interpreter resource counts above are independently measured. Model-supplied certificates are reproduced as claims and are not automatically promoted to proved results.

\lstset{literate={−}{{$-$}}1 {²}{{$^2$}}1 {³}{{$^3$}}1 {¹}{{$^1$}}1 {⁴}{{$^4$}}1 {⁵}{{$^5$}}1 {⁷}{{$^7$}}1 {⁹}{{$^9$}}1 {₀}{{$_0$}}1 {₁}{{$_1$}}1 {₂}{{$_2$}}1 {≤}{{$\leq$}}1 {±}{{$\pm$}}1 {·}{{$\cdot$}}1 {×}{{$\times$}}1 {→}{{$\to$}}1 {–}{{--}}1}

\paragraph{Square request.}

\lstinputlisting[basicstyle=\ttfamily\scriptsize,breaklines=true,breakatwhitespace=false,columns=fullflexible,keepspaces=true]{architecture_square_164_prompt.txt}

\paragraph{Square complete final answer.}

\lstinputlisting[basicstyle=\ttfamily\scriptsize,breaklines=true,breakatwhitespace=false,columns=fullflexible,keepspaces=true]{architecture_square_164_reply.txt}

\paragraph{Sine request.}

\lstinputlisting[basicstyle=\ttfamily\scriptsize,breaklines=true,breakatwhitespace=false,columns=fullflexible,keepspaces=true]{architecture_sin_164_prompt.txt}

\paragraph{Sine complete final answer.}

\lstinputlisting[basicstyle=\ttfamily\scriptsize,breaklines=true,breakatwhitespace=false,columns=fullflexible,keepspaces=true]{architecture_sin_164_reply.txt}

\subsection{All public construction summaries}\label{app:architecture_summaries}

Each paragraph reproduces the inferred rule, implementation explanation and certificate fields returned by one request. Unicode mathematical characters are typeset equivalently. Failed executable graphs are labeled below. Their explanations remain part of the response record.

\paragraph{Square, seed 164, valid.}

\textbf{Family.} Single multiplicative gate with no stored coefficients


\textbf{Inferred rule.} y = x\textsuperscript{2}


\textbf{Construction summary.} Multiply the scalar input by itself. The graph uses one node, zero stored coefficients, one scalar multiplication, and one output scalar, with no intermediate activations. It implements the inferred square rule throughout [-10,10], including beyond the observed input range.


\textbf{Model-supplied certificate.} For every real x, the graph computes x\ensuremath{\cdot}x = x\textsuperscript{2}. The observed pairs agree with this rule to binary64 precision. Binary64 inference performs one rounded multiplication; on [-10,10] there is no overflow, and absolute rounding error is at most 2\textasciicircum{}-47. Finite examples do not uniquely determine a rule, so this certificate establishes implementation of the inferred rule rather than uniqueness of inference.


\paragraph{Square, seed 317, valid.}

\textbf{Family.} Parameter-free multiplicative graph


\textbf{Inferred rule.} y = x\textsuperscript{2}


\textbf{Construction summary.} A single multiplication node uses the scalar input for both operands. It stores zero numerical coefficients, performs one scalar multiplication, and produces one scalar output without intermediate activations. The construction implements the inferred rule throughout [-10,10], including beyond the observed inputs.


\textbf{Model-supplied certificate.} In exact arithmetic, the node returns x\ensuremath{\cdot}x = x\textsuperscript{2} for every real x. In binary64 it returns the rounded product, with no overflow on [-10,10]; underflow is possible only for extremely small inputs. The observed outputs agree with the square rule to displayed numerical precision. Finite examples do not uniquely determine a function, so this certificate establishes implementation of the inferred rule, not uniqueness of the inference.


\paragraph{Square, seed 583, valid.}

\textbf{Family.} Parameter-free multiplicative network


\textbf{Inferred rule.} y = x\textsuperscript{2}


\textbf{Construction summary.} A single multiplication node uses the scalar input as both operands. The network stores zero numerical coefficients and performs one scalar multiplication, producing one output scalar with no intermediate activations. It implements the inferred square rule throughout [-10,10], including outside the observed input range.


\textbf{Model-supplied certificate.} In real arithmetic, the node outputs input \ensuremath{\times} input = x\textsuperscript{2} identically. In binary64, it outputs the correctly rounded product fl(x \ensuremath{\times} x); for |x| \ensuremath{\leq} 10, no overflow occurs and absolute rounding error is at most 2\textasciicircum{}(-47). The observed pairs are consistent with the square rule at the displayed precision. Finite examples alone do not uniquely determine the underlying function.


\paragraph{Sine, seed 164, valid.}

\textbf{Family.} Scalar polynomial gating network with triple-angle reconstruction


\textbf{Inferred rule.} The examples agree with y = sin(x), inferred over the full domain [-10,10]. Finite examples alone do not uniquely determine this rule.


\textbf{Construction summary.} Scale the input by 1/27, evaluate the degree-13 sine Taylor polynomial using Horner gates, and apply s \ensuremath{\to} s(3-4s\textsuperscript{2}) three times. This implements sine across the deployment domain with near-binary64 accuracy, without a sine operation. The graph uses 23 nodes and 19 stored scalar coefficients. With activation reuse, inference requires only a few live scalars and 34 scalar arithmetic operations. No global resource optimality is claimed.


\textbf{Model-supplied certificate.} Let P(t)=t-t\textsuperscript{3}/3!+t\textsuperscript{5}/5!-t\textsuperscript{7}/7!+t\textsuperscript{9}/9!-t\textsuperscript{1}\textsuperscript{1}/11!+t\textsuperscript{1}\textsuperscript{3}/13! and F(s)=3s-4s\textsuperscript{3}. For |t|\ensuremath{\leq}10/27, the alternating-series bound gives |P(t)-sin(t)|\ensuremath{\leq}(10/27)\textasciicircum{}15/15!<2.7e-19. Also F composed three times at sin(t) equals sin(27t). On the intervening initial-value interval, the derivative of this composition has magnitude below 30, using sin(27 arcsin(s)); hence the exact-coefficient real-arithmetic error is below 8.1e-18 throughout [-10,10]. Stored-coefficient rounding and binary64 arithmetic add floating-point error; the real-arithmetic bound is not a bound on that additional error.


\paragraph{Sine, seed 317, valid.}

\textbf{Family.} Scalar polynomial network with triple-angle gates


\textbf{Inferred rule.} The examples are consistent with y = sin(x), with x in radians. Finite examples alone do not uniquely determine this rule.


\textbf{Construction summary.} Scale the input by 1/27 and evaluate the degree-13 sine Taylor polynomial using Horner evaluation. Three triple-angle gates then restore the original argument. This implements the inferred sine rule throughout [-10,10], rather than interpolating the displayed samples. The graph stores 14 scalar coefficients. Sequential evaluation with activation reuse requires only a few live scalars. No global resource optimality is claimed.


\textbf{Model-supplied certificate.} Let p(t)=t-t\textasciicircum{}3/3!+t\textasciicircum{}5/5!-t\textasciicircum{}7/7!+t\textasciicircum{}9/9!-t\textasciicircum{}11/11!+t\textasciicircum{}13/13!, and T(s)=3s-4s\textasciicircum{}3. In exact arithmetic the graph computes T(T(T(p(x/27)))). For |t|<=10/27, the alternating-series bound gives |p(t)-sin(t)|<2.6e-19. On this range arcsin has derivative below 1.08, so |arcsin(p(t))-t|<2.808e-19. Since T(sin(a))=sin(3a), the uniform exact-arithmetic output error is below 7.6e-18. Stored coefficient rounding and binary64 arithmetic add floating-point roundoff; the exact-arithmetic bound is not a claim of correctly rounded sine.


\paragraph{Sine, seed 583, valid.}

\textbf{Family.} Polynomial gating network with shared triple-angle coefficients


\textbf{Inferred rule.} y = sin(x), consistent with every displayed pair.


\textbf{Construction summary.} Scale the input by 1/27, evaluate the degree-13 sine Taylor polynomial using Horner's rule, and apply s \ensuremath{\to} s(3 - 4s\textsuperscript{2}) three times. This reconstructs sine across the full deployment interval without a sine operation. The graph stores 10 scalar coefficients and uses 27 scalar arithmetic operations. Sequential evaluation needs only a small constant workspace. This is a compact rule-based construction, not a claim of global optimality.


\textbf{Model-supplied certificate.} Let t=x/27 and P(t)=sum from k=0 to 6 of (-1)\textasciicircum{}k t\textasciicircum{}(2k+1)/(2k+1)!. For |x|\ensuremath{\leq}10, |P(t)-sin(t)|\ensuremath{\leq}(10/27)\textasciicircum{}15/15!<2.7e-19. The identity 3sin(u)-4sin(u)\textasciicircum{}3=sin(3u) makes the output sin(27 arcsin(P(t))). Since |sin(t)|<0.363, the mean-value theorem bounds its ideal real-arithmetic error by 27(2.7e-19)/sqrt(1-0.364\textsuperscript{2})<8e-18 uniformly. Stored decimal coefficients and binary64 operations introduce additional rounding error; the bound above describes the approximation error, not a certified floating-point bound. Finite observations cannot uniquely determine the generating rule.


\paragraph{Positive self-power, seed 164, valid.}

\textbf{Family.} Rule-based scalar network with polynomial logarithm and scaled polynomial exponential


\textbf{Inferred rule.} The examples are consistent with y = x\textasciicircum{}x. The network approximates this rule throughout 0 < x <= 10.


\textbf{Construction summary.} The network computes x log(x) using a reciprocal iteration followed by a 64-term odd-power logarithm series. It then approximates the exponential with a fourth-degree Taylor polynomial and eight squarings. All coefficients are explicit. The logarithm approximation is most accurate away from zero; its absolute contribution to the exponent still vanishes as x approaches zero. This is a constructive approximation with an explicit error certificate, not a claim that finite examples uniquely identify the rule or that the architecture is globally optimal.


\textbf{Model-supplied certificate.} In exact arithmetic, the reciprocal after seven iterations is r=(1-e)/(x+1), where e=((5-x)/6)\textasciicircum{}128 and 0<=e<7.4e-11. Put z=(x-1)r and a=|z|. The computed exponent is L=2x sum\_\{k=0\}\textasciicircum{}\{63\} z\textasciicircum{}(2k+1)/(2k+1). Its error from x log(x) is at most D=7.4e-10+2x*a\textasciicircum{}129/(129*(1-a*a)). The output is P4(L/256)\textasciicircum{}256, where P4(t)=1+t+t\textasciicircum{}2/2+t\textasciicircum{}3/6+t\textasciicircum{}4/24. Since |L|<=24, define E=exp(2*24/256)*(24/256)\textasciicircum{}5/120. Relative output error is at most exp(D+256*E/(1-E))-1. These bounds exclude binary64 rounding. The logarithm tail bound tends to zero as x approaches zero when its actual series remainder, rather than the geometric upper bound, is used.


\paragraph{Positive self-power, seed 317, invalid.}

\textbf{Family.} Scalar gated arithmetic network with rational quadrature and polynomial exponentiation


\textbf{Inferred rule.} y = x\textasciicircum{}x. The graph approximates this rule over 0 < x <= 10.


\textbf{Construction summary.} Approximate log(x) by 16-point Gauss--Legendre quadrature of (x-1)/(1+t(x-1)). Fourteen parallel Newton iterations compute the quadrature reciprocals. Multiply by x, evaluate a fourth-degree exponential polynomial at one 256th of the result, and square eight times. This construction follows a domain-wide mathematical rule rather than interpolating the observations. It uses 64 nodes and 59 stored scalar coefficients. Accuracy is strongest away from zero; the certificate supplies a conservative uniform relative-error bound, not a claim of machine-precision accuracy or global architectural optimality.


\textbf{Model-supplied certificate.} The observations are consistent with exp(x log(x)). For exact reciprocals, quadrature is exact through degree 31. Expanding the integral geometrically bounds the log error by (1-x)\textasciicircum{}32/33 for x<=1 and (1-1/x)\textasciicircum{}32/33 for x>=1. Thus the error in x log(x) is at most 1/33 throughout the domain. Each normalized reciprocal has residual (1-d)\textasciicircum{}16384; d>=0.0005299532504175, so its residual is below 0.000170. For x<=1 its additional exponent error is below 0.000170/e; for x>=1 it is negligible. The scaled exponent lies between -1/(256e) and 10 log(10)/256. Taylor's remainder bounds the final exponentiation relative error below 0.000014. Together these give a conservative uniform relative error below 0.032, with ample margin for binary64 rounding. The graph is an approximation, not an exact representation of x\textasciicircum{}x.


\paragraph{Positive self-power, seed 583, invalid.}

\textbf{Family.} Scalar arithmetic network with quadrature, shared vector operations, and repeated squaring


\textbf{Inferred rule.} y = x\textasciicircum{}x. The network approximates exp(x log x) on 0 < x <= 10, with limit 1 as x approaches zero.


\textbf{Construction summary.} A 32-point Gauss--Legendre quadrature approximates log(x) using an integral of reciprocal affine functions. Thirteen vectorized Newton steps approximate those reciprocals. Multiplication by x produces the exponent. A fourth-order exponential polynomial followed by twelve squarings produces the output. This is a rule-based approximation rather than an interpolation of the observations. Accuracy is strongest away from zero; the small-input region has additional quadrature and reciprocal approximation error. No global resource optimality is claimed.


\textbf{Model-supplied certificate.} The observations agree with x\textasciicircum{}x. Use log(x)=(x-1) integral\_0\textasciicircum{}1 1/(1+(x-1)t) dt. The listed t,w are 32-point Gauss--Legendre nodes and weights. Write d\_i=1+(x-1)t\_i. Since 0<d\_i<10, Newton iteration r'=r(2-d\_i r), initialized at 0.1, gives r\_i=(1-(1-0.1d\_i)\textasciicircum{}8192)/d\_i after 13 steps. Thus the graph's exponent is exactly z=x(x-1) sum\_i w\_i r\_i in real arithmetic. Its output is [1+u+u\textsuperscript{2}/2+u\textsuperscript{3}/6+u\textsuperscript{4}/24]\textasciicircum{}4096, u=z/4096. These identities provide a directly checkable approximation certificate, not an assertion of exact equality to x\textasciicircum{}x. Binary64 adds rounding error.


\paragraph{Hyperbolic tangent, seed 164, valid.}

\textbf{Family.} Single scalar tanh activation with no learned coefficients


\textbf{Inferred rule.} y = tanh(x).


\textbf{Construction summary.} Apply the standard tanh activation directly to the scalar input. This uses one graph node, zero stored numerical coefficients, and one scalar activation evaluation. The same rule applies throughout [-10,10], including the saturation regions beyond the observed examples.


\textbf{Model-supplied certificate.} The observed pairs agree with y = tanh(x) to displayed numerical precision. The graph computes this inferred rule directly for every input in [-10,10], subject only to binary64 activation rounding. Finite samples cannot uniquely determine an arbitrary generating function; this certificate establishes implementation of the inferred rule, not uniqueness of the inference.


\paragraph{Hyperbolic tangent, seed 317, valid.}

\textbf{Family.} Parameter-free scalar tanh network


\textbf{Inferred rule.} y = tanh(x).


\textbf{Construction summary.} Apply one tanh activation directly to the scalar input. The network stores zero numerical coefficients, uses one scalar activation output, and requires one tanh evaluation with no additional arithmetic. It implements the inferred rule across the entire deployment interval [-10,10].


\textbf{Model-supplied certificate.} The sole node computes tanh(x) by definition, so in exact arithmetic the graph equals the inferred rule for every real x. Binary64 inference introduces only the activation implementation's numerical error. The observed pairs agree with this rule to the displayed precision; finite samples alone do not uniquely determine an underlying function.


\paragraph{Hyperbolic tangent, seed 583, valid.}

\textbf{Family.} Single scalar tanh activation with no stored coefficients


\textbf{Inferred rule.} y = tanh(x).


\textbf{Construction summary.} Apply the standard tanh activation directly to the scalar input. The network uses one node, zero stored coefficients, and one scalar activation evaluation, with no affine arithmetic or intermediate layers. It implements the inferred rule throughout [-10,10], including the saturation regions beyond the observed inputs.


\textbf{Model-supplied certificate.} The graph directly evaluates tanh(x), so its output equals the inferred rule over the entire deployment domain, subject only to binary64 activation rounding. The displayed pairs agree with this rule to their shown numerical precision. Finite observations do not uniquely determine an underlying function; this certificate establishes implementation of the inferred rule rather than uniqueness of inference.


\paragraph{Addition, seed 164, valid.}

\textbf{Family.} Parameter-free additive network


\textbf{Inferred rule.} y = x\textsubscript{1} + x\textsubscript{2}.


\textbf{Construction summary.} Extract the two input coordinates and add them. The graph stores no numerical coefficients and performs one binary64 addition, with no hidden layers or nonlinearities. It implements the inferred rule throughout the deployment domain.


\textbf{Model-supplied certificate.} For every input (a,b) in [-10,10]\textsuperscript{2}, the take nodes return [a] and [b], so the output is [fl(a+b)], where fl denotes binary64 rounding. The observed outputs agree with this rule to displayed precision. Finite observations do not uniquely determine an unrestricted underlying function.


\paragraph{Addition, seed 317, valid.}

\textbf{Family.} Parameter-free additive graph


\textbf{Inferred rule.} y = x\textsubscript{0} + x\textsubscript{1}.


\textbf{Construction summary.} Select the two input coordinates and add them. This implementation stores no numerical coefficients and performs one binary64 addition. It requires no hidden layers or nonlinear activations and applies throughout the deployment domain.


\textbf{Model-supplied certificate.} The take nodes return [x\textsubscript{0}] and [x\textsubscript{1}], so the output is [fl(x\textsubscript{0} + x\textsubscript{1})], where fl denotes binary64 rounding. The observed outputs agree with this rule to floating-point precision. The same identity holds for every input in [-10,10]\textsuperscript{2}; the examples support, but do not uniquely determine, the inferred rule.


\paragraph{Addition, seed 583, valid.}

\textbf{Family.} Parameter-free additive graph


\textbf{Inferred rule.} y = x\textsubscript{0} + x\textsubscript{1}


\textbf{Construction summary.} Extract the two input coordinates and add them directly. The graph stores no numerical weights, biases, or constants and performs one scalar addition. It uses two one-entry selections and a one-entry output. This construction implements the inferred rule throughout [-10,10]\textsuperscript{2}, including inputs far beyond the observed examples.


\textbf{Model-supplied certificate.} Every displayed output agrees with x\textsubscript{0} + x\textsubscript{1} to binary64 rounding precision. For any input, the take nodes return x\textsubscript{0} and x\textsubscript{1} exactly, so the output is fl(x\textsubscript{0} + x\textsubscript{1}), the correctly rounded binary64 sum. Thus the graph implements the inferred addition rule over the full domain. Finite observations alone do not uniquely determine the underlying rule.


\paragraph{Subtraction, seed 164, valid.}

\textbf{Family.} Single linear neuron


\textbf{Inferred rule.} Output the first input coordinate minus the second: y = x\textsubscript{1} - x\textsubscript{2}.


\textbf{Construction summary.} One bias-free affine node with weights [1, -1] computes the difference. It stores two scalar coefficients and produces one activation scalar, with no hidden layers or nonlinearities. The affine expression requires two multiplications and one addition; multiplication by \ensuremath{\pm}1 can be implemented as identity and sign change.


\textbf{Model-supplied certificate.} For every input (a,b) in [-10,10]\textsuperscript{2}, [1,-1]\ensuremath{\cdot}[a,b] = a-b. This matches the displayed outputs to binary64 rounding precision and implements the inferred rule throughout the domain. Multiplication by \ensuremath{\pm}1 is exact in binary64, leaving only the rounding of the subtraction.


\paragraph{Subtraction, seed 317, valid.}

\textbf{Family.} Single bias-free linear neuron


\textbf{Inferred rule.} y = x\textsubscript{0} - x\textsubscript{1}


\textbf{Construction summary.} A single affine node uses weights [1, -1] and no bias. It stores two scalar coefficients and produces one output activation, with no hidden layers or intermediate activation tensors. The same subtraction rule applies throughout the deployment domain.


\textbf{Model-supplied certificate.} For every input, [1, -1] \ensuremath{\cdot} [x\textsubscript{0}, x\textsubscript{1}] = x\textsubscript{0} - x\textsubscript{1}. This matches all displayed outputs to binary64 rounding precision. Multiplication by \ensuremath{\pm}1 is exact on the specified domain, so evaluation yields the correctly rounded subtraction throughout [-10,10]\textsuperscript{2}. The examples support this rule but cannot uniquely determine an unrestricted function.


\paragraph{Subtraction, seed 583, valid.}

\textbf{Family.} Single-layer linear network


\textbf{Inferred rule.} y = x\textsubscript{1} - x\textsubscript{2}.


\textbf{Construction summary.} One affine node subtracts the second coordinate from the first using weights [1, -1] and no bias. It stores two scalar coefficients and produces one activation scalar, with no hidden layers. The same rule applies throughout [-10,10]\textsuperscript{2}.


\textbf{Model-supplied certificate.} For every input (a,b), the network computes [1,-1]\ensuremath{\cdot}[a,b] = a-b in exact arithmetic. All displayed outputs agree with this rule to binary64 rounding precision. The weights are exactly representable in binary64. This certifies implementation of the inferred rule across the domain; finite observations alone cannot uniquely determine an arbitrary unseen target.


\paragraph{Running sum, seed 164, valid.}

\textbf{Family.} Single-state additive recurrent network


\textbf{Inferred rule.} Each output is the running sum of the inputs through that position.


\textbf{Construction summary.} Maintain one scalar state initialized to zero. Add each incoming scalar to the state and emit the updated state directly. The network stores one scalar coefficient, uses one addition per element, and requires constant scalar workspace apart from input/output storage. No learned weights or nonlinearities are needed, and the same cell handles arbitrary sequence lengths.


\textbf{Model-supplied certificate.} In exact arithmetic, s\_0=0 and s\_t=s\_\{t-1\}+x\_t imply by induction that s\_t=sum\_\{i=1\}\textasciicircum{}t x\_i. Thus the output is the prefix sum for every sequence over [-10,10], regardless of length. In binary64, the cell implements sequential rounded summation, matching the observed rule up to floating-point rounding.


\paragraph{Running sum, seed 317, valid.}

\textbf{Family.} Scalar linear recurrent accumulator


\textbf{Inferred rule.} Each output is the inclusive prefix sum: y\_t = x\_1 + ... + x\_t.


\textbf{Construction summary.} Use one scalar state initialized to zero. Each cell adds the current input to the state, and the readout returns the updated state directly. The network stores one coefficient (the initial zero), performs one addition per input, and needs constant workspace independent of sequence length, excluding retained outputs. It applies to arbitrary sequence lengths and inputs throughout [-10,10], without clipping or saturation.


\textbf{Model-supplied certificate.} In exact arithmetic, s\_0 = 0 and s\_t = s\_\{t-1\} + x\_t imply by induction that s\_t = sum\_\{i=1\}\textasciicircum{}t x\_i. Returning s\_t therefore gives every inclusive prefix sum. In binary64, the same cell implements sequential floating-point prefix summation, with its ordinary rounding and representational limits.


\paragraph{Running sum, seed 583, valid.}

\textbf{Family.} Single-state additive recurrent network


\textbf{Inferred rule.} Each output is the inclusive prefix sum of the input sequence: y\_t = x\_1 + ... + x\_t.


\textbf{Construction summary.} Initialize one scalar state to zero. At each step, add the current input to the state and emit the updated state directly. The network stores one scalar coefficient, uses one addition per element, and requires constant streaming workspace, excluding retained inputs and outputs. No trained weights or nonlinearities are needed. The same recurrence applies to arbitrary sequence lengths and all inputs in [-10,10].


\textbf{Model-supplied certificate.} Initially s\_0 = 0. If s\_(t-1) equals the sum of the first t-1 inputs, the cell gives s\_t = s\_(t-1) + x\_t, the sum of the first t inputs. Induction proves the prefix-sum rule in exact arithmetic. Binary64 inference computes the corresponding left-to-right floating-point sums, subject to rounding and representable range.


\paragraph{Sort four, seed 164, valid.}

\textbf{Family.} Sparse ReLU sorting network with five compare-exchange units


\textbf{Inferred rule.} Return the four input coordinates in ascending order.


\textbf{Construction summary.} Three stages compare two disjoint pairs, compare their respective minima and maxima, then order the two middle candidates. Each compare-exchange uses subtraction, ReLU, and addition, implemented with one shared stored coefficient, -1. Indexing and concatenation route values without dense matrices. The construction implements sorting throughout [-10,10]\textasciicircum{}4, including ties and unseen permutations, subject to binary64 rounding. No global resource optimality is claimed.


\textbf{Model-supplied certificate.} For a pair (a,b), let t=ReLU(a-b). Then (a-t,b+t)=(min(a,b),max(a,b)) in exact arithmetic. After the first stage, write the sorted pairs as l0<=h0 and l1<=h1. The second stage produces L=min(l0,l1), U=min(h0,h1), V=max(l0,l1), H=max(h0,h1). L is the global minimum, H the global maximum, and U,V are the remaining elements. Sorting U,V therefore sorts all four values. Binary64 introduces only rounding error; on the stated domain a conservative absolute error bound per output relative to exact sorting is 1e-13.


\paragraph{Sort four, seed 317, valid.}

\textbf{Family.} Three-stage ReLU sorting network


\textbf{Inferred rule.} Return the four input coordinates in ascending order.


\textbf{Construction summary.} Five compare-exchange units sort four values in three stages. Each unit computes r = ReLU(a - b), then returns (a - r, b + r). The first stage sorts two input pairs; the second compares their minima and their maxima; the final stage sorts the middle two values. Parallel units share convolution coefficients, and all stages share the correction vector. The graph stores eight scalar coefficients, including zeros and ones, and uses no biases. It implements sorting throughout the deployment domain, with only binary64 rounding error.


\textbf{Model-supplied certificate.} In exact arithmetic, (a-ReLU(a-b), b+ReLU(a-b)) = (min(a,b),max(a,b)). After stage one, write the pairs as l0\ensuremath{\leq}h0 and l1\ensuremath{\leq}h1. Stage two produces [min(l0,l1),max(l0,l1),min(h0,h1),max(h0,h1)]. Its first and last entries are the global minimum and maximum. Sorting its two middle entries therefore sorts all four coordinates. This proof includes ties and holds for all real inputs. Binary64 subtraction and addition can introduce rounding; exact bitwise preservation is not asserted.


\paragraph{Sort four, seed 583, valid.}

\textbf{Family.} Five-comparator ReLU sorting network with one shared scalar coefficient


\textbf{Inferred rule.} Return the four input coordinates in ascending order.


\textbf{Construction summary.} Three stages perform five compare-and-exchange operations: two on the original pairs, two on their minima and maxima, and one on the middle candidates. Each comparator computes r = ReLU(a - b), then returns (a - r, b + r). Vector operations execute independent comparators together. The only stored numerical coefficient is a shared -1; indexing and concatenation specify the wiring. Arithmetic comprises 15 scalar additions, 10 scalar negations implemented as multiplications, and 5 ReLUs. Intermediate tensors have at most four entries. No fitting or domain-specific thresholds are used.


\textbf{Model-supplied certificate.} In real arithmetic, r=max(a-b,0) gives a-r=min(a,b) and b+r=max(a,b). After stage one, let the sorted pairs be (p,P) and (q,Q). Stage two produces L=min(p,q), u=min(P,Q), v=max(p,q), H=max(P,Q). L is the global minimum, H is the global maximum, and the remaining values are u,v. Sorting u,v therefore sorts all four inputs, including ties. This identity holds for every real input. Binary64 evaluation introduces ordinary subtraction/addition roundoff; the stated domain cannot overflow.


\paragraph{Fibonacci, seed 164, valid.}

\textbf{Family.} Two-state additive recurrent network


\textbf{Inferred rule.} Fibonacci sequence: F(0)=0, F(1)=1, and F(n+2)=F(n+1)+F(n).


\textbf{Construction summary.} Initialize the state to [0,1]. Each iteration replaces [a,b] with [b,a+b]. After n iterations, return the first coordinate. The network stores only two numerical initialization coefficients and uses one scalar addition per iteration, with constant-size inference workspace. Indexing and concatenation require no scalar arithmetic. This implements the inferred recurrence across the entire deployment domain rather than approximating the displayed values.


\textbf{Model-supplied certificate.} Induction gives state=[F(n),F(n+1)] after n iterations: initialization establishes n=0, and the cell implements the defining recurrence. For n<=64, every state entry and addition result is a nonnegative integer at most F(65)=17167680177565<2\textasciicircum{}53, so all arithmetic is exact in binary64.


\paragraph{Fibonacci, seed 317, valid.}

\textbf{Family.} Two-state additive recurrent network


\textbf{Inferred rule.} Fibonacci sequence: F(0)=0, F(1)=1, and F(n+2)=F(n)+F(n+1).


\textbf{Construction summary.} Initialize two state entries to [0,1]. Each iteration replaces [a,b] with [b,a+b]. After n iterations, return the first entry. The network stores only two numerical coefficients, uses one scalar addition per iteration, and requires constant workspace. Selection and concatenation require no scalar arithmetic. This construction implements the inferred recurrence throughout the deployment domain rather than interpolating the examples.


\textbf{Model-supplied certificate.} Induction gives state=[F(n),F(n+1)] after n iterations: initialization establishes the base case, and the update establishes the induction step. For n<=64, every computed integer is at most F(65)=17167680177565<2\textasciicircum{}53, so all additions and outputs are exact in binary64. The observed examples agree with this rule, although finitely many examples cannot uniquely determine an extrapolation.


\paragraph{Fibonacci, seed 583, valid.}

\textbf{Family.} Two-state additive recurrent network


\textbf{Inferred rule.} Fibonacci sequence: F(0)=0, F(1)=1, and F(n+2)=F(n+1)+F(n).


\textbf{Construction summary.} Initialize the state to [0,1]. Each iteration replaces [a,b] with [b,a+b]. After n iterations, return the first entry. The construction stores two numerical coefficients in the initial state and uses one scalar addition per iteration, with constant-size workspace and no learned weights. It implements the inferred recurrence throughout the deployment domain rather than interpolating the examples.


\textbf{Model-supplied certificate.} The initial state is [F(0),F(1)]. If the state is [F(k),F(k+1)], the cell produces [F(k+1),F(k)+F(k+1)]=[F(k+1),F(k+2)]. Induction proves the output is F(n). For n<=64, every state entry is a nonnegative integer at most F(65)=17167680177565<2\textasciicircum{}53, so binary64 additions and outputs are exact.


\paragraph{Triangular, seed 164, valid.}

\textbf{Family.} Scalar quadratic network with multiplicative gating


\textbf{Inferred rule.} The triangular-number continuation: y(n) = n(n + 1)/2.


\textbf{Construction summary.} Square the input, add the original input, and scale by one half. This graph stores one scalar coefficient and uses three scalar arithmetic operations with constant-size workspace. It implements the inferred rule across the entire deployment domain without iteration or a lookup table.


\textbf{Model-supplied certificate.} The observed successive differences are 1,2,...,8, consistent with y(n)=sum\_\{k=1\}\textasciicircum{}n k. The graph computes (n\textsuperscript{2}+n)/2, which equals that sum for every nonnegative integer n. For 0\ensuremath{\leq}n\ensuremath{\leq}64, n\textsuperscript{2}\ensuremath{\leq}4096 and n\textsuperscript{2}+n\ensuremath{\leq}4160; these integers and multiplication by 0.5 are represented exactly in binary64. The finite observations do not uniquely determine the continuation.


\paragraph{Triangular, seed 317, valid.}

\textbf{Family.} Scalar polynomial network with multiplicative gating


\textbf{Inferred rule.} Triangular numbers: y(n) = n(n + 1)/2.


\textbf{Construction summary.} An affine node computes (n + 1)/2, then a multiplicative gate multiplies it by n. This construction uses two nodes, two stored scalar coefficients, and three scalar arithmetic operations per input, with constant workspace and no iteration.


\textbf{Model-supplied certificate.} The output is n(0.5n + 0.5) = n(n + 1)/2. It is zero at n=0 and its successive difference is n, matching all observed pairs. For every integer 0 \ensuremath{\leq} n \ensuremath{\leq} 64, all intermediate values and outputs are exactly representable in binary64, so evaluation is exact. The examples support this rule but do not uniquely determine an arbitrary unseen function.


\paragraph{Triangular, seed 583, valid.}

\textbf{Family.} Scalar affine network with multiplicative gating


\textbf{Inferred rule.} Triangular numbers: y(n) = n(n + 1)/2.


\textbf{Construction summary.} An affine node computes (n + 1)/2, and a multiplicative gate multiplies it by n. The network stores two scalar coefficients and uses two scalar multiplications and one addition per inference, with constant workspace and no iteration.


\textbf{Model-supplied certificate.} The output is n(0.5n + 0.5) = n(n + 1)/2. This gives y(0) = 0 and y(n + 1) - y(n) = n + 1, matching the observed successive increments. For every integer 0 <= n <= 64, all intermediate values are exactly representable in binary64, so the inferred rule is computed exactly. Finite examples do not uniquely determine extrapolation.

